\chapter{一変数の積分}
\label{ch:tb-integral}

\section{出題で使う積分操作}

変数上端には微積分学の基本定理
\[
\frac{\d}{\d x}\int_a^x f(t)\,\d t=f(x)
\]
を使う。置換積分と部分積分は
\[
\int_{\phi(\alpha)}^{\phi(\beta)}f(x)\,\d x
=\int_\alpha^\beta f(\phi(t))\phi'(t)\,\d t,
\]
\[
\int_a^bu(x)v'(x)\,\d x
=\bigl[u(x)v(x)\bigr]_a^b-\int_a^bu'(x)v(x)\,\d x
\]
である。定積分の置換では端点も変える。

2015年度第1期で必要な部分積分は、$a\ne-1$ のとき
\[
\int x^a\log x\,\d x
=x^{a+1}\left(\frac{\log x}{a+1}-\frac1{(a+1)^2}\right)+C.
\]
$a=-1$ だけは $u=\log x$ と置いて
\[
\int\frac{\log x}{x}\,\d x=\frac{(\log x)^2}{2}+C
\]
となる。$\sqrt{x^2+1}$ の置換は第\ref{ch:tb-functions}章に置いた。

\section{広義積分と比較}

非有界区間または特異点を含む積分は、端を切った積分の極限で定義する。
例えば
\[
\int_a^\infty f(x)\,\d x
=\lim_{R\to\infty}\int_a^Rf(x)\,\d x.
\]
特異点が内部にあれば左右を別々に判定する。

\[
\int_1^\infty x^a\,\d x<\infty\Longleftrightarrow a<-1,
\qquad
\int_0^1x^a\,\d x<\infty\Longleftrightarrow a>-1.
\]
$0\le f\le g$ かつ $\int g<\infty$ なら $\int f<\infty$ であり、
$\int\abs f<\infty$ なら $\int f$ は収束する。

問題の $\int_1^\infty x^a\log x\,\d x$ は $t=\log x$ により
\[
\int_0^\infty te^{(a+1)t}\,\d t
\]
へ変わるため、$a<-1$ のとき、かつそのときに限り収束する。

\section{局所化する指数核}

\begin{theorem}{近似恒等核 $ne^{-nx}$}{tb-exponential-kernel}
$f:[0,\infty)\to\R$ が有界かつ $0$ で連続なら
\[
\lim_{n\to\infty}\int_0^\infty n f(x)e^{-nx}\,\d x=f(0).
\]
\end{theorem}

\begin{proof}
$\int_0^\infty ne^{-nx}\,\d x=1$ なので、差は
\[
\int_0^\infty n\bigl(f(x)-f(0)\bigr)e^{-nx}\,\d x
\]
である。$\varepsilon>0$ に対し、$0\le x\le\delta$ で
$\abs{f(x)-f(0)}<\varepsilon$ となる $\delta$ を取る。
$M=\sup_{x\ge0}\abs{f(x)-f(0)}$ とすれば、積分を $\delta$ で分けて
\[
\left|\int_0^\infty n(f(x)-f(0))e^{-nx}\,\d x\right|
\le\varepsilon+Me^{-n\delta}\longrightarrow\varepsilon.
\]
$\varepsilon$ は任意なので結論を得る。
\end{proof}
